\begin{annexitem}{annex:p12.har.asymmetry}{thm:p12-har-asymmetry}{Every application in the relation proof}
\begin{proof}
In the context \(X:\mathrm{Set},R:\mathrm{Rel}\,X,
h:\mathrm{AntiSym}\,X R,x:X,q:Rxx\), unfolding \(\mathrm{Rel}\)
gives \(R:X\to X\to\mathrm{Prop}\), so \(Rxx:\mathrm{Prop}\).
Unfolding \(\mathrm{AntiSym}\) gives
\(h:\Pi u,v:X.Ruv\to Rvu\to\bot\). Application to \(x:X\)
gives \(hx:\Pi v:X.Rxv\to Rvx\to\bot\); the second application
gives \(hxx:Rxx\to Rxx\to\bot\). Applying to \(q:Rxx\)
gives \(hxxq:Rxx\to\bot\), and applying again gives
\(hxxqq:\bot\). Abstraction over \(q\) has type
\(Rxx\to\bot\); abstraction over \(x\) has type
\(\Pi x:X.Rxx\to\bot\), convertible to \(\mathrm{Irrefl}\,X R\).
Abstracting over \(h,R,X\), in that order, gives precisely the
claimed type. The \(R,X\) binders are dependent because later
annotations refer to them. Domain comparisons are \(X\) against \(X\)
and \(Rxx\) against \(Rxx\); definition unfolding and beta contraction
are the only changes in type. This verifies the whole example, not
synthesis for every term of the source system.
\end{proof}
\end{annexitem}

\begin{annexitem}{annex:p12.har.normal-form-test}{thm:p12-har-normal-form-test}{Finite search and both directions of conversion}
\begin{proof}
Enumerate the finite reduction tree breadth first, retaining each finite
path. Effective finite branching makes each level computable. Test each
encountered term with the decidable normal-form predicate. By hypothesis
some finite path reaches a normal form, so the search halts. Type
preservation keeps each test within the specified well-typed regime.
This procedure does not assume that every arbitrary reduction strategy
terminates.

Let \(m,n\) be computed normal forms of \(M,N\). If they are
alpha-equivalent, their reduction paths exhibit conversion. Conversely,
confluence gives the Church--Rosser property for a finite conversion
zigzag. Prove this by induction along the zigzag, maintaining a common
reduct of its endpoints. When the new edge is forward, join its reduct
with the earlier common reduct by confluence; when it is backward, the
new endpoint reduces to the previous one, hence to the earlier common
reduct. This proves the induction step. Apply it to the zigzag from
\(m\) back to \(M\), across to \(N\), and down to \(n\).
Since \(m,n\) are normal, their common reduct can only be themselves
up to alpha-equivalence. Comparing these normal forms therefore decides
conversion.

For a system with sound and complete synthesis and unique types up to this
conversion, a synthesized type \(C\) gives \(M:A\) exactly when
\(C=A\): conversion proves one direction, uniqueness proves the other.
If synthesis reports no type, completeness excludes a derivation. These
premises are not consequences of the generic term grammar.
\end{proof}
\end{annexitem}

\begin{annexitem}{annex:p12.har.framework-forward}{thm:p12-har-framework-forward}{All object-rule cases for the three signatures}
\begin{proof}
All substitutions are capture avoiding; alpha-rename bound variables
outside the finite free-variable sets of inserted terms. In the
propositional signature, translate each hypothesis \(A\) to a proof
variable \(h:T[A]\). Induct on a natural-deduction derivation, keeping
the undischarged hypothesis list. The hypothesis case is its variable.
For implication introduction, the induction term \(u:T[B]\) under
\(h:T[A]\) gives \(\lambda h:T[A].u:T[A]\to T[B]\).
Applying \(\mathrm{imp\_intr}[A][B]\) gives
\(T(\mathrm{imp}[A][B])=T[A\supset B]\); the discharged
variable is bound. For elimination, induction terms
\(u:T[A\supset B]\), \(v:T[A]\) give
\(\mathrm{imp\_el}[A][B]uv:T[B]\). These are all logical
rules of this fragment. Weakening adds an unused declaration, exchange
leaves binding unchanged, and repeated use of a hypothesis uses the same
variable again. The host context discipline permits these structural actions.

For predicate translation, induction on object terms gives host terms of
type \(A\): the cases are variables, \(c\), and application of
\(f\). Induction on formulas gives names of type \(\mathrm{prop}\):
atomic \(Rst\), implication, and
\(\mathrm{forall}(\lambda x:A.[\varphi])\).
The same structural inductions prove
\([\varphi[t/x]]=[\varphi][[t]/x]\) up to host beta/alpha
conversion. The variable case performs precisely the designated
substitution; a different variable or constant is unchanged; application
commutes with the induction hypotheses; in the binder case a fresh
binder prevents capture and the hypothesis applies to its body.

Hypothesis and implication derivation cases are as above. In universal
introduction, the eigenvariable \(x\) occurs in none of the
undischarged assumptions. The induction term \(u:T[\varphi]\)
under a fresh \(x:A\) gives
\(\lambda x:A.u:\Pi x:A.T[\varphi]\). Applying
\(\mathrm{forall\_intr}(\lambda x:A.[\varphi])\) gives
\(T[\forall x:A.\varphi]\), with no illicit dependence of an
undischarged assumption on \(x\). Universal elimination applies
\(\mathrm{forall\_elim}(\lambda x:A.[\varphi])u[t]\), with
type \(T([\varphi][[t]/x])\), convertible to
\(T[\varphi[t/x]]\). These exhaust the predicate fragment.
In the source example, the first universal elimination on \(h:T U\)
at \(z:A\) gives \(T(\mathrm{forall}(\lambda y:A.Rzy))\);
the second at \(z\) gives \(T(Rzz)\). Implication introduction
closes \(h\), and universal introduction closes \(z\), exactly
as in the explicit body term.

For untyped terms, structural induction shows their translations have
type \(D\): variables are declared, application uses \(\mathrm{app}\),
and abstraction applies \(\mathrm{abs}\) to a host function
\(D\to D\). The same variable, application and fresh-binder
cases prove \([P[Q/x]]=[P][[Q]/x]\) up to host conversion.
Induct on beta-equational derivations. Reflexivity uses
\(\mathrm{refl}[P]\), symmetry uses
\(\mathrm{sym}[P][Q]u\), and transitivity uses
\(\mathrm{trans}[P][Q][R]uv\). Application congruence uses
\(\mathrm{mon}\) with the argument equality first and function
equality second, matching its source type. Abstraction congruence
abstracts the induction proof over fresh \(x:D\), then applies
\(\mathrm{xi}\) to the two translated host functions. For the beta
generator put \(F=\lambda x:D.[P]\). The term
\(\mathrm{beta}\,F[Q]\) has type
\(\mathrm{eq}(\mathrm{app}(\mathrm{abs}\,F)[Q])(F[Q])\).
Its first argument is \([(\lambda x.P)Q]\); its second
beta-reduces to \([P][[Q]/x]\), hence to \([P[Q/x]]\).
Alpha renaming is absorbed by the host binding convention. These exhaust
the beta-equality generators and congruence cases; no eta generator is added.

The constructions start from object derivations, not arbitrary host terms,
so they prove no converse faithfulness claim.
\end{proof}
\end{annexitem}

\begin{annexitem}{annex:p12.har.sigma-extraction}{thm:p12-har-sigma-extraction}{Dependent projection types and reduction}
\begin{proof}
For \(x:A\), application yields \(px:\Sigma y:B.Rxy\).
First projection gives \(\pi_1(px):B\), hence abstraction gives
\(f:A\to B\). Second projection gives
\(\pi_2(px):Rx(\pi_1(px))\), with its dependency on the actual
first projection. The defined \(fx\) beta/delta-reduces to
\(\pi_1(px)\), so conversion gives \(\pi_2(px):Rx(fx)\),
and abstraction gives the claimed \(q\).

If \(p\twoheadrightarrow\lambda x:A.\langle t,u\rangle\),
application to the bound argument in \(f\) reduces to
\(\langle t,u\rangle\), with capture-avoiding substitution;
projection computation then gives \(t\). Thus the function reduces
to \(\lambda x:A.t\). Second projection computes the associated
proof in the same way. All these terms use the admitted sum eliminators.
If existential elimination into data is forbidden, the first projection
is not licensed for that existential and this construction cannot be used.
\end{proof}
\end{annexitem}

\begin{annexitem}{annex:p12.har.reflection}{thm:p12-har-reflection}{The two conversion premises}
\begin{proof}
Application gives
\(\mathrm{ok}\,t:(\mathrm{Check}\,t)\to(\mathrm{Dec}\,t)\).
The checked conversion \(\mathrm{Check}\,t=_{\beta\iota\delta}\top\)
lets \(\mathrm{true}:\top\) be used at type
\(\mathrm{Check}\,t\). Application yields
\(\mathrm{ok}\,t\,\mathrm{true}:\mathrm{Dec}\,t\).
The other checked conversion changes this type to \(T\). Neither
checker completeness nor an externally trusted success answer was used.
\end{proof}
\end{annexitem}

\begin{annexitem}{annex:p12.har.leibniz}{thm:p12-har-leibniz}{Equivalence and transport as explicit terms}
\begin{proof}
Write \(\mathrm{Eq}\,a b\) for Leibniz equality. Reflexivity is
\(r_a:=\lambda P:A\to\mathrm{Prop}.\lambda u:Pa.u\), of type
\(\Pi P:A\to\mathrm{Prop}.Pa\to Pa\).
For \(e:\mathrm{Eq}\,a b\), choose the predicate
\(P_0:=\lambda z:A.\mathrm{Eq}\,z a\). It has the allowed
proposition-valued type in this impredicative regime. Since
\(r_a:P_0a\), the term \(eP_0r_a\) has type
\(P_0b=\mathrm{Eq}\,b a\). Thus symmetry was constructed,
not assumed in defining the predicate.

For \(e:\mathrm{Eq}\,a b\), \(d:\mathrm{Eq}\,b c\),
the term \(\lambda P:A\to\mathrm{Prop}.\lambda u:Pa.dP(ePu)\)
has type \(\mathrm{Eq}\,a c\), proving transitivity.
For any admitted \(P:A\to\mathrm{Prop}\) and \(u:Pa\), the
definition directly gives \(ePu:Pb\). No data-valued predicate was
inserted into the permitted proposition-valued class.
\end{proof}
\end{annexitem}

\begin{annexitem}{annex:p12.har.setoid-constructions}{thm:p12-har-setoid-constructions}{Function equality, subset equality, and quotient maps}
\begin{proof}
Write \(f=\langle f_0,f_1\rangle\),
\(g=\langle g_0,g_1\rangle\). At each argument, reflexivity,
symmetry and transitivity of \(=_B\) respectively give those three
properties for pointwise function equality. These operations never
inspect the preservation proofs, so equal underlying functions with
arbitrary preservation proofs give equal function-space values.

Given pointwise equality \(e\) and \(d:x=_A y\), preservation
by \(g\) gives \(g_0x=_B g_0y\). Composing with
\(e x:f_0x=_B g_0x\) gives \(f_0x=_B g_0y\).
Conversely, the assertion at all equal pairs specializes to \(y=x\)
with the reflexivity proof of \(=_A\), giving pointwise equality.

For subsetoid elements \(\langle a,u\rangle,\langle b,v\rangle\),
the declared equality is just \(a=_A b\); its three equivalence
operations give the three operations on the subset carrier without
consulting \(u,v\). Both elements already belong to the carrier,
so no implication \(Pa\to Pb\) or predicate invariance was assumed.

The quotient carrier remains \(A\) with equality \(Q\), assumed
an equivalence relation. A term \(f:A\to B\) with
\(\Pi x,y:A.Qxy\to fx=_B fy\) therefore has exactly the proof
required for a map from \([A,Q]\). Precomposing that proof with
\(=_A\subseteq Q\) also proves preservation of the old equality.
No representative selector, choice axiom or changed underlying function
is needed.
\end{proof}
\end{annexitem}

\begin{annexitem}{annex:p12.har.field-independence}{thm:p12-har-field-independence}{Reciprocal proofs and the four absolute-value branch pairs}
\begin{proof}
Subsetoid equality identifies \(\langle a,p\rangle\) and
\(\langle a',q\rangle\) exactly when \(a=_A a'\), independently
of \(p,q\). The reciprocal's preservation proof gives equal outputs.
For a change of proof alone, take \(a'=a\) and reflexivity.

For absolute value assume \(x=_F x'\). Transport branch inequalities
along this equality, using the order's assumed respect for it. If both
branches are nonnegative, the outputs \(x,x'\) are equal. If both
are nonpositive, their outputs \(-x,-x'\) are equal because negation
respects equality. If the first branch is nonnegative and the second
nonpositive, the transported inequalities give \(x\ge0\) and
\(x\le0\). Order antisymmetry gives \(x=_F0\), hence also
\(x'=_F0\). Using \(-0=_F0\) and equality-preserving negation
shows \(x=_F-x'\). The reversed mixed pair gives \(-x=_F x'\)
by the same two inequalities. These four cases exhaust the sum constructors,
including their overlap at zero. The argument requires the ordered-field
laws; it does not derive them from the sign-proof type or identify proofs
definitionally.
\end{proof}
\end{annexitem}

\begin{annexitem}{annex:p12.har.hol-forward}{thm:p12-har-hol-forward}{Every direct HOL translation clause}
\begin{proof}
First induct on domains. A basic domain has its declared Type sort;
\(\Omega\) translates to \(\mathrm{Prop}:\mathrm{Type}\).
If the two domain translations have sort Type, their arrow has sort Type
by the allowed \((\mathrm{Type},\mathrm{Type})\) pair. They have
no dependence on the arrow's fresh variable.

Induct on object-term formation. Constants and variables use their
declarations and weakening into the complete signature. Application uses
the product/application rule at the translated domain arrow. Abstraction
uses that same well-formed arrow and the body induction judgment.
For logical implication the two formula induction results have sort Prop,
so \((\mathrm{Prop},\mathrm{Prop})\) forms the translated
implication. Universal quantification has Type-sorted domain and
Prop-sorted body, so \((\mathrm{Type},\mathrm{Prop})\) forms
the translated quantifier. These cases establish the term/formula claim.

Induction on the term grammar proves substitution commutation. A
designated variable is replaced on both sides, other variables/constants
stay unchanged, application uses the two hypotheses, and a fresh binder
allows the hypothesis to act on its body without capture. Consequently
translation sends a beta contraction to a beta contraction with the
same substituted body. It therefore preserves finite beta conversion.

Now induct on the HOL derivation. An assumption is its proof variable
\(z_i:[\varphi_i]\). Implication introduction abstracts the induction
term over the designated proof variable, with the Prop/Prop product;
elimination applies the term for the implication to the term for its
antecedent. Universal introduction abstracts over \(x:[\sigma]\)
with the Type/Prop product. Its eigenvariable condition ensures that
none of the undischarged proof-variable types depends on \(x\), so
this binder can be discharged while those assumptions remain in context.
Universal elimination applies the proof term to \([t]\) and yields
the type \([\varphi][[t]/x]=[\varphi[t/x]]\).
For the HOL conversion rule, the previous structural calculation gives
\([\varphi]=_\beta[\psi]\); the translated \([\psi]\)
already has sort Prop, so the lambdaHOL conversion rule gives the new
type with unchanged proof term. These exhaust Figure 1.
\end{proof}
\end{annexitem}

\begin{annexitem}{annex:p12.har.connective-checks}{thm:p12-har-connective-checks}{All encoded connective cases and source symmetry}
\begin{proof}
Expand abbreviations before reduction; no delta rule is added to the
formal lambdaHOL calculus. For conjunction, a term
\(p:\Pi\alpha:\mathrm{Prop}.(\varphi\to\psi\to\alpha)\to\alpha\)
applied to \(\varphi:\mathrm{Prop}\) expects a term of
\(\varphi\to\psi\to\varphi\). The first projection supplies
\(\lambda x:\varphi.\lambda y:\psi.x\), so its result is
\(\varphi\). At \(\psi\), the second projection supplies
\(\lambda x:\varphi.\lambda y:\psi.y\), so its result is
\(\psi\). For \(u:\varphi,v:\psi\), abstraction over
\(\alpha:\mathrm{Prop}\) and
\(h:\varphi\to\psi\to\alpha\) gives \(huv:\alpha\),
hence the encoded pair has the conjunction type. Every function arrow
inside is formed by Prop/Prop; quantification over \(\alpha\) uses
Type/Prop since \(\mathrm{Prop}:\mathrm{Type}\).

After expanding the first projection of that pair, beta reduction gives
\[
 (\lambda\alpha:\mathrm{Prop}.
  \lambda h:\varphi\to\psi\to\alpha.huv)
       \varphi(\lambda x:\varphi.\lambda y:\psi.x)
 \twoheadrightarrow_\beta
 (\lambda x:\varphi.\lambda y:\psi.x)uv
 \twoheadrightarrow_\beta u.
\]
Replacing its last body by \(y\) and its first argument by \(\psi\)
gives the second computation \(v\).

For disjunction, introductions from \(u:\varphi\), \(v:\psi\)
are respectively
\[
 \lambda\alpha:\mathrm{Prop}.
 \lambda f:\varphi\to\alpha.\lambda g:\psi\to\alpha.fu,
 \qquad
 \lambda\alpha:\mathrm{Prop}.
 \lambda f:\varphi\to\alpha.\lambda g:\psi\to\alpha.gv.
\]
Both inhabit the stated disjunction. An encoded disjunction proof
applied to \(\chi\) and its two implication proofs has type
\(\chi\). Substituting either introduction into this elimination
beta-reduces to \(fu\) or \(gv\), respectively. A falsehood proof
has type \(\Pi\alpha:\mathrm{Prop}.\alpha\), so application to
\(\chi:\mathrm{Prop}\) gives \(\chi\), as required.

For the existential, given \(t:A\) and
\(u:\varphi[t/x]\), a hypothetical
\(h:\Pi x:A.\varphi\to\alpha\) applied to \(t,u\) gives
\(\alpha\). Abstracting over \(h,\alpha\) gives the encoded
existential proof. Conversely its application to \(\chi\) and
\(h:\Pi x:A.\varphi\to\chi\) yields \(\chi\).
The requirement that \(x\) not occur in \(\chi\) or the
undischarged external assumptions is precisely what lets the local
witness/hypothesis binders construct \(h\). Eliminating the displayed
introduction beta-reduces to \(htu\). This proves the local rules,
not extraction to a Type-sorted data conclusion.

Finally let \(r:x=_A y\), \(P:A\to\mathrm{Prop}\).
The predicate \(\lambda z:A.Pz\to Px\) has type
\(A\to\mathrm{Prop}\); applying \(r\) to it gives
\((Px\to Px)\to(Py\to Px)\), after beta conversion.
The term \(\lambda q:Px.q\) supplies its first input, yielding
\(Py\to Px\). Abstracting over \(P\) gives \(y=_A x\),
and then over \(r\) gives the exact source symmetry type.
\end{proof}
\end{annexitem}

\begin{annexitem}{annex:p12.har.hol-decoder}{thm:p12-har-hol-decoder}{Inverse decoding with its explicit structural premises}
\begin{proof}
Use \cref{thm:p12-har-disambiguation} to place the context in its three strata.
Decode a simple data type by reading a declared domain variable as its
HOL basic domain, Prop as \(\Omega\), and arrows as domain arrows.
Decode object variables/constants, application and abstraction structurally.
The Type/Type restriction supplies exactly the domain arrow formation;
no proof declaration is needed for these terms. For a Prop-sorted formula,
a product with Prop-sorted domain is nondependent by disambiguation, and
decodes to implication; a product with Type-sorted domain decodes to
universal quantification. The remaining object-term/application/abstraction
cases decode by the object grammar, including proposition-valued predicates.
These operations commute with substitution by the same variable,
application and fresh-binder induction as the forward translation.
In particular beta conversion of admitted formulas decodes to HOL beta
conversion.

Induct on the supplied proof typing derivation, using generation to
identify its underlying type-forming case beneath conversion/weakening.
For a variable of proposition type, the context-strata premise places it
in \(\Gamma_P\), so it decodes to that assumption. A sort axiom
cannot have a proposition type: its right side is Type or Type'. A
product expression itself has a sort rather than a proposition as its
type, so it cannot be the subject of this proof case either.

An abstraction proving a product of sort Prop has a domain of sort
Prop or Type. In the first case its body type does not depend on the
proof binder; apply the body induction construction and implication
introduction, discharging its formula assumption. In the second case
decode its data binder, apply the body construction, and use universal
introduction. Freshness of the original context binder ensures that it
is absent from all earlier external assumptions, giving the eigenvariable
condition. These are precisely the two allowed product pairs with
propositional codomain.

In an application whose result is a proposition type, generation exposes
the function's product type. Type correctness and the allowed pairs again
make its domain Prop- or Type-sorted. For a Prop domain, decode both
subderivation proofs and apply implication elimination. For a Type domain,
decode the function proof and the argument as an object term; universal
elimination gives the decoded substituted body. Substitution commutation
identifies it with the required decoded result formula. A Type-sorted
codomain product cannot produce this proof case, by the assumed
type/sort correctness and classification.

Weakening adds a declaration: if it is a proof declaration, add its
formula as an unused assumption; if it is a domain/object declaration,
extend the HOL signature. For conversion, first decode the premise,
then apply HOL conversion to the two beta-equal admitted formulas.
Context exchange is justified by \cref{thm:p12-har-disambiguation}.
This exhausts the typing rules. Generation, classification and type
correctness are supplied by
\cref{thm:p12-har-hol-generation,thm:p12-har-disambiguation,thm:p12-har-hol-type-correctness};
this induction constructs the inverse direct HOL derivation using them.
It neither assumes that inverse as a premise nor applies it to an
unrelated logical-framework signature.
\end{proof}
\end{annexitem}

\begin{annexitem}{annex:p12.har.guard-counterexample}{thm:p12-har-guard-counterexample}{The missing Type/Prop branch}
\begin{proof}
The context is well formed: \(A:\mathrm{Type}\) is licensed by
\(\mathrm{Type}:\mathrm{Type}'\); \(A\to\mathrm{Prop}\)
has sort Type by Type/Type; and \(\Pi x:A.Px\) has sort Prop
by Type/Prop. In the extension by a fresh \(x:A\), application
gives \(hx:Px\). The same product formation then permits
\(\lambda x:A.hx:\Pi x:A.Px\).

The body synthesizes to \(Px\), an application, not a simple arrow
ending in Prop or a domain variable. It is also neither Type nor Type'.
Thus the printed modified guard selects its second branch and tests
whether \(\mathrm{Type}_\Gamma(A)=\mathrm{Prop}\). Variable
lookup instead returns Type, so the guard rejects the abstraction.
The corrected second branch admits either Prop or Type and therefore
accepts this valid Type/Prop case. This is a concrete failure of the
printed completeness assertion, not a counterexample to the source's
Figure 2 calculus.
\end{proof}
\end{annexitem}

\begin{annexitem}{annex:p12.har.checker-conditional}{thm:p12-har-checker-conditional}{The repaired guard, mutual soundness, termination and completeness}
\begin{proof}
The structural facts used below are proved in
\cref{thm:p12-har-hol-contexts,thm:p12-har-hol-generation,thm:p12-har-hol-type-correctness,thm:p12-har-hol-context-conversion,thm:p12-har-disambiguation,thm:p12-har-hol-uniqueness,thm:p12-har-hol-subject-reduction},
with raw confluence and substitution in
\cref{thm:p12-har-pseudo-confluence,thm:p12-har-pts-substitution}.
Strong normalization is proved independently in
\cref{thm:p12-har-hol-normalization}; no checker or termination
claim is used in that normalization argument.
First justify the guard on its stated premise: a successfully typed
body \(M:B\) in a well-formed extension. Type correctness gives
either a sort for \(B\) or an exceptional untyped top-sort case.
The disambiguation classification identifies a data-type shape exactly
with \(B:\mathrm{Type}\). A product with that codomain is
admitted exactly for a Type-sorted domain, since Type/Type is the only
allowed pair ending in Type. If \(B\) is neither data-shaped nor
Type nor Type', the same classification gives \(B:\mathrm{Prop}\).
The product is then admitted exactly when the domain has sort Prop or
Type. If \(B\) is Type, its sort is Type', with no allowed product
pair ending in Type'; if it is Type', no formation sort exists. These
last cases must fail. This proves the corrected guard's equivalence to
the product-formation premise for this particular typed-body regime.

Prove partial soundness mutually by induction on finite evaluation
trees of Type and Ok. Empty context is valid. An accepted extension has
a fresh variable and a type whose recursively synthesized sort is
sound, so the variable/weakening rules make the context valid. Variable
lookup retrieves its context type; the accepted-context sort branches
are the axioms weakened into that context. The top sort never succeeds.
The product branch's two successful recursive judgments and allowed
pair give exactly the product rule. In abstraction, the recursive body
judgment is sound, and the guard equivalence just established supplies
the required well-formed product; hence abstraction is sound.

For application, recursive soundness gives \(M:C\), \(N:D\).
Type correctness handles their returned types. Subject reduction and
generation keep the weak-head exposed product \(\Pi x:A.B\)
well formed at its required sort; in particular its domain has an
admitted sort. Conversion of \(M\) to this product type and of
\(N:D\) to the beta-equal domain \(A\) permits application.
Its output is exactly \(B[N/x]\). This case uses the stated
metatheoretic premises, rather than trusting the comparison result
without a typed conversion judgment.

For termination choose explicit weights: a variable or sort has weight
one, application has the sum of its children plus one, and product or
abstraction has the sum of annotation and body plus two. A context
extension has the old context's weight plus the annotation's weight plus
one. The input weight is the context weight for Ok and the sum of
context and term weights for Type. Every recursive call is smaller:
Ok on an extension calls Type on its annotation in the shorter context;
a variable/sort calls Ok on its context; product and application call
their proper subterms, with one context-extension unit instead of the
two binder units for the bound body; abstraction calls its bound body
and its annotation in precisely that way. The original recursive call
on a newly computed product has been removed, so it is not counted as
a decreasing call.

Induct on this weight. All recursive children terminate. By the partial
soundness just proved, successful application children return well-typed
types. Their reductions terminate by \cref{thm:p12-har-hol-normalization};
normal-form equality is decidable by confluence, and weak-head product
exposure therefore terminates as well. Other tests are finite syntactic
or context-membership tests. This proves total termination even on
untypable inputs: no conversion test is run before its recursive typed
inputs have succeeded. The original generic pseudo-term equality problem
was never claimed decidable.

For completeness, use generation to strip the possible final conversion
and weakening, and induct on the resulting syntax-directed typing
premises together with well-formed context extension. Context formation
requires the annotation's admitted sort, which the induction result
synthesizes; sorts do not beta-convert to a different sort. Therefore Ok
accepts exactly those extensions. Variables and the two typed sorts use
the accepted context and return their prescribed types; no judgment
types Type'. In a product judgment, the two induction results return
its domain/body sorts and the admitted pair is recognized. In an
abstraction judgment, the body result is beta-equal to its derivable
body type by uniqueness; subject reduction, conversion, type correctness
and disambiguation classify it at the same permitted codomain sort.
The guard consequently accepts exactly the required pair, and the
returned product is beta-equal to the derivable abstraction type.

In an application judgment the induction results give a synthesized
function type \(C\) beta-equal to a derivable
\(\Pi x:A.B\), and an argument type beta-equal to \(A\).
Confluence and normalization make weak-head normalization of \(C\)
expose a product with beta-equal domain and codomain, rather than a
nonproduct head. Domain comparison succeeds. Substitution preserves
beta conversion, so its returned codomain instance is beta-equal to
the derivable result type. Conversion at the outer judgment preserves
that equality, and weakening changes no synthesized meaning in the
accepted extended context. The given generation/context-conversion
premises justify these rearrangements; their proofs are not an output
of the checker itself.

Thus synthesis succeeds precisely on typable terms and yields a type
beta-equal to any derivable type. Soundness, uniqueness and decidable
conversion give the final TCP comparison in both directions. The
normalization theorem is used through its separate proof, not inferred
from this checker's asserted termination.
\end{proof}
\end{annexitem}

\begin{annexitem}{annex:p12.har.pts-morphism}{thm:p12-har-pts-morphism}{All PTS rule cases and the normalization transfer}
\begin{proof}
Induct on pseudo terms to prove
\(f(M[N/x])=f(M)[f(N)/x]\): variables/sorts give the direct
clauses; application uses its two induction hypotheses; abstraction
and product use a fresh binder and their annotation/body hypotheses.
Thus a beta contraction maps to the beta contraction of its mapped
abstraction with the mapped argument. A contraction in an annotation
or body is preserved by the homomorphic compatible contexts. Finite
beta conversion is consequently preserved as well.

Induct on a source typing derivation. A sort axiom is mapped to an
admitted target axiom by the defining preservation condition. Variable
and weakening use the recursively mapped annotation judgment; fresh
variables remain fresh because \(f\) fixes variables. In product
formation the two mapped sort premises combine with the preserved
triple. Abstraction combines its mapped body and mapped product-sort
premise. Application combines its mapped function and argument and
uses substitution commutation to identify the result type. Conversion
uses the mapped well-formed result type and preserved beta equality.
These are all seven general PTS rules.

For the displayed PRED-omega/HOL map, the two sort axioms become
Prop:Type and Type:Type'. Its six product pairs map respectively into
Type/Type, Type/Type, Prop/Prop, Type/Prop, Type/Type and Type/Prop,
all admitted in lambdaHOL. For the map to CC, both axioms become
\(\ast:\Box\), and its pairs become among CC's four allowed pairs.
These checks establish the actual morphisms, not a name-based system
identification.

Finally suppose a well-typed source term had an infinite beta-contraction
sequence. Map every designated contraction at its redex position.
The result is an infinite sequence of target contractions starting
from the well-typed image term, contradicting target strong
normalization. Sort identification does not erase the abstraction and
application constructors of a redex. This argument needs only preservation
of contraction steps, not an inverse relation on arbitrary preimages.
\end{proof}
\end{annexitem}

\begin{annexitem}{annex:p12.har.pseudo-confluence}{thm:p12-har-pseudo-confluence}{Parallel beta reduction with annotations and dependent products}
\begin{proof}
Use an auxiliary reflexive parallel relation, denoted here by
\(\leadsto\), on the exact source pseudo-term grammar. A variable
or sort relates to itself. Abstractions and products relate by relating
their annotations and bodies under the same fresh binder. Applications
relate by relating their two children. In addition, if
\(A\leadsto A'\), \(M\leadsto M'\),
\(N\leadsto N'\), the redex
\((\lambda x:A.M)N\) relates to \(M'[N'/x]\).
These are a proof device for beta confluence, not new typing rules or an
extension of the source conversion relation.

Induction on the term grammar gives reflexivity. Induction on a
parallel-reduction derivation proves substitution compatibility:
if \(M\leadsto M'\) and \(N\leadsto N'\), then
\(M[N/x]\leadsto M'[N'/x]\). For a variable equal to \(x\)
this is the second premise; for another variable or a sort it is
reflexivity. Application uses the two child inductions. In abstraction
and product choose the binder fresh for \(N,N'\), apply the
inductions to annotation and body, and use the compatible binder case.
For a parallel contraction, the two substitutions commute after a
fresh binder choice, and the induction hypotheses apply to its body
and argument, giving the substituted contraction case. Thus annotations
and dependent product bodies have not been ignored.

Define a complete development \(T^*\) recursively. Variables/sorts
are unchanged; an abstraction or product develops its annotation/body;
an application \(MN\) with literal abstraction
\(M=\lambda x:A.P\) develops to \(P^*[N^*/x]\); any other
application develops to \(M^*N^*\). This finite structural definition
does not require normalization and need not contract redexes newly
created by a substitution.

Prove by structural induction on the original term that
\(T\leadsto U\) implies \(U\leadsto T^*\).
For sorts/variables there is only reflexivity. At an abstraction or
product, its child parallel steps use the induction hypotheses and
compatible binder step. For a non-redex application, the original
parallel derivation must be its compatible application case, giving
\(M'\leadsto M^*\), \(N'\leadsto N^*\); compatible application
then gives \(M'N'\leadsto M^*N^*\). This remains valid even if
\(M'\) has become an abstraction: choose the compatible application
case rather than a root contraction. For an original literal redex,
there are two possibilities. If its parallel step preserves the
application constructor, it has result
\((\lambda x:A'.P')N'\); the child inductions give parallel
steps to \(A^*,P^*,N^*\), and a root parallel contraction gives
\(P^*[N^*/x]\). If its parallel step contracts the root, its
result is \(P'[N'/x]\); substitution compatibility with the same
child inductions gives that result parallel to \(P^*[N^*/x]\).
These exhaust the two application cases and prove the absorption property.

Consequently, if \(T\leadsto U,V\), both \(U,V\) relate in
one parallel step to \(T^*\); parallel reduction is diamond.
A finite rectangular grid of these diamonds joins any two finite
parallel paths: fill the first row against the first edge of the second
path, then successive rows, using a diamond at every cell. The bottom
and right edges yield the common endpoint. Reflexive edges cover empty
paths. Thus the reflexive-transitive parallel relation is confluent.

Every single compatible beta contraction is parallel: relate unaffected
children reflexively, perform the designated root parallel contraction,
and lift it through its term context. Conversely every parallel step
is a finite compatible beta path. Prove this by induction on its
derivation: the compatible cases reduce each child in turn, and in a
contraction case first reduce annotation, body and argument by those
induction paths, then contract the resulting root redex. Therefore the
reflexive-transitive closures of beta and parallel reduction coincide.
Confluence of the latter is exactly the required confluence of beta.
No typed preservation or strong-normalization premise entered the proof.
\end{proof}
\end{annexitem}

\begin{annexitem}{annex:p12.har.pts-substitution}{thm:p12-har-pts-substitution}{Every source PTS rule under context substitution}
\begin{proof}
First record the free-variable discipline by induction on Figure 3
derivations: every free variable of the subject and its type is declared
in the context. Sort axioms have none; variable/weakening use the
annotation premise; application combines its two judgments; abstraction
and product bind their new variable; conversion uses its separate
well-typed result-type premise. Thus expressions typed in
\(\Gamma_1\) contain no free \(x\), and the given \(N\)
has no free variable declared after \(x\). Freshen later declaration
names and binders outside the free variables of \(N\).

Induct on the source derivation, considering its chosen split at \(x\).
For a premise context that precedes the declaration of \(x\),
substitution has no effect by that discipline; retain its derivation.
For a premise containing \(x\), use the induction hypothesis with
the inherited suffix split. These two alternatives cover premises on
both sides of the declaration.

A sort axiom has empty context, so is in the no-effect case. For a
variable conclusion, let its final declaration be \(y:C\).
If \(y=x\), the suffix is empty, its subject is \(x\), and
its annotation is the selected \(B\). The substituted conclusion
is exactly the given \(\Gamma_1\vdash N:B\).
If \(y\ne x\) and the selected declaration is in its prefix,
apply the induction hypothesis to the annotation judgment
\(C:s\). It yields \(C[N/x]:s\) in the substituted prefix.
The variable rule adds fresh \(y:C[N/x]\) and concludes the
required substituted variable judgment. If the entire variable
conclusion precedes \(x\), retain it unchanged as already described.

For weakening by \(y:C\), if \(y=x\) and it is the selected
declaration, both subject/type in the weakening premise lie in
\(\Gamma_1\). Substitution leaves them unchanged; removing this
one final weakening step gives the desired judgment in \(\Gamma_1\).
Otherwise substitute its annotation and subject premises by induction,
or retain a premise preceding \(x\). Weakening by
\(y:C[N/x]\) then gives the substituted conclusion. In particular,
an old occurrence of \(x\) carried through later weakenings becomes
the given \(N\) and is carried through these new well-formed later
declarations; no unproved context-thinning rule is needed.

In product formation, the bound \(y\) is fresh for \(x,N\).
The induction judgments give the substituted domain at its original
sort \(s_1\) and the substituted codomain in the context extended
by \(y:C[N/x]\) at \(s_2\). The same specification triple
\((s_1,s_2,s_3)\) therefore forms the substituted product at
\(s_3\). For abstraction, substitute the body judgment in that
fresh extended context and the product-sort premise in the base
context, then use the abstraction rule. These are exactly the
capture-avoiding substitutions in annotation and body.

For application, induction gives the substituted function and argument
judgments. The result of their application is the original codomain
with first \(x\) and then the fresh bound-variable argument
substitution. Capture-avoiding substitution composition identifies it
with the old result type \(A\) substituted by \(N/x\).
This composition identity follows structurally: variables give the two
designated replacements, application combines them, and fresh binders
allow the same identity on annotations and bodies.

For conversion, induction gives the substituted subject premise and
the substituted well-typed result type at its same sort. Beta equality
is preserved by substitution: a substituted beta generator contracts
to the substituted contractum after the same fresh-binder composition
identity, and compatible contexts/equivalence steps are preserved.
The conversion rule therefore changes to the required substituted type.
These cases exhaust Figure 3, including later context types that
depend on \(x\). The construction produces a derivation in the
substituted context, which consequently is well formed in the source's
sense whenever it occurs. No normalization or strengthening was used.
\end{proof}
\end{annexitem}

\begin{annexitem}{annex:p12.har.inverse-counterexample}{thm:p12-har-inverse-counterexample}{The missing mixed predicate-domain/data-codomain pair}
\begin{proof}
In lambdaHOL, \(A:\mathrm{Type}\) is a valid context.
\(A\) and Prop both have sort Type, so Type/Type forms
\(A\to\mathrm{Prop}:\mathrm{Type}\). Using Type/Type again,
with \(A:\mathrm{Type}\) as codomain, forms
\((A\to\mathrm{Prop})\to A:\mathrm{Type}\). The bound
variables of both arrows are unused.

The printed reverse sends the domain declaration to \(A:\mathrm{Set}\).
In PRED-omega, the axiom Prop:Type and pair Set/Type give
\(A\to\mathrm{Prop}:\mathrm{Type}\); the context declaration
gives \(A:\mathrm{Set}\). Therefore an outer product with
that literal domain and codomain requires the pair Type/Set.
The exact source specification has only Set/Set, Set/Type,
Type/Type, Prop/Prop, Set/Prop and Type/Prop. None is Type/Set.

Conversion does not repair this missing pair. Types of the literal
context variable \(A\) are beta-equal to its annotation Set:
induct on its variable, weakening and conversion derivation cases.
No other subject-constructor rule can conclude a variable judgment.
Similarly, the inner literal product's sort must be Type by its only
applicable Set/Type formation pair, up to conversion. The sort constants
Set and Type are distinct beta-normal forms, hence not beta-equal by
the proved raw confluence theorem. Stripping final conversions and
weakenings from a putative outer-product derivation must therefore still
reach the absent Type/Set formation case. No such derivation exists.

Since the source reverse classifies an arrow chain ending in the
variable \(A\) as a Set-valued type, it applies precisely to this
judgment and fails. This refutes that proposed reverse map, not every
possible encoding of HOL into another formal system.
\end{proof}
\end{annexitem}
\begin{annexitem}{annex:p12.har.hol-contexts}{thm:p12-har-hol-contexts}{Valid contexts, projection, and thinning}
\begin{proof}
Define a valid context inductively: the empty context is valid; a fresh
extension \(\Gamma,x:A\) is valid when \(\Gamma\) is valid and
\(\Gamma\vdash A:s\) for a source sort \(s\). Every context in a
source derivation is valid. This follows by induction on that derivation:
the sort cases have empty context; variable and weakening exhibit the
annotation certificate; product and abstraction inherit their base
context from a premise; application and conversion inherit their common
context. Conversely every valid context admits
\(\Gamma\vdash\mathrm{Prop}:\mathrm{Type}\): start from the axiom and
weaken through its annotation certificates. Thus this explicit validity
is equivalent to the source's occurrence-in-a-judgment definition.

If \(x:A\) occurs in a valid context, the variable rule types \(x:A\)
at its own prefix and successive weakenings carry that judgment through
the suffix. This proves projection without a new primitive rule.
If \(\Gamma\) and \(\Delta\) are valid, every declaration of
\(\Gamma\) occurs literally in \(\Delta\), and
\(\Gamma\vdash M:A\), then \(\Delta\vdash M:A\).
Induct on the displayed derivation. For a sort, weaken its empty-context
axiom through \(\Delta\). For a variable, use the projection just proved.
A weakening case uses the induction hypothesis on its subject premise;
its formerly added declaration is already present in \(\Delta\).
Product and abstraction use the induction hypotheses on their premises,
renaming the bound variable fresh for \(\Delta\); the new annotation
judgment makes \(\Delta,x:B\) valid, so the body induction applies
there. Application combines the two induction judgments. Conversion
combines its subject and well-typed result-type induction judgments and
the unchanged raw beta equality. This covers all Figure~2 rules.
In particular this thinning proof does not infer strengthening by
substituting a possibly nonexistent inhabitant.
\end{proof}
\end{annexitem}

\begin{annexitem}{annex:p12.har.hol-generation}{thm:p12-har-hol-generation}{Generation and product injectivity}
\begin{proof}
For every source judgment, generation has the following exact content.
For a variable \(x:T\) in the context, any type of \(x\) is
beta-equal to \(T\). Any type of \(\mathrm{Prop}\) is beta-equal to
\(\mathrm{Type}\); any type of \(\mathrm{Type}\) is beta-equal to
\(\mathrm{Type}'\); no judgment has subject \(\mathrm{Type}'\).
For a product \(\Pi x:B.C\), there are premises
\(\Gamma\vdash B:s_1\),
\(\Gamma,x:B\vdash C:s_2\), an allowed pair, and its conclusion type
is beta-equal to \(s_2\). For an abstraction \(\lambda x:B.P\),
there is a \(C\) with \(\Gamma,x:B\vdash P:C\),
\(\Gamma\vdash\Pi x:B.C:s\), and conclusion type beta-equal to
\(\Pi x:B.C\). For an application \(PQ\), there are
\(\Gamma\vdash P:\Pi x:B.C\), \(\Gamma\vdash Q:B\), and its
conclusion type is beta-equal to \(C[Q/x]\).
The displayed constructor judgment at this generated type is itself
derivable in \(\Gamma\).

Prove these statements by induction on the typing derivation. Its
constructor rule supplies the witnesses. A final conversion composes
the generated type equality with its own equality. A final weakening
carries every generated premise into the enlarged valid context by \cref{thm:p12-har-hol-contexts},
including the fresh binder extension, and leaves the equalities intact.
No other rule has the selected subject constructor. These are all cases.

Raw beta confluence gives product injectivity without normalization.
If \(\Pi x:B.C=_\beta\Pi x:D.E\), choose a common reduct.
A product has no root beta contraction, and every reduct still has
the product constructor with a reduced domain and body. After choosing
the same fresh binder, the common reduct therefore identifies common
reducts of \(B,D\) and of \(C,E\). Hence
\(B=_\beta D\), \(C=_\beta E\). Distinct sorts are not beta-equal,
and a product is not beta-equal to a sort, by the same argument.
\end{proof}
\end{annexitem}

\begin{annexitem}{annex:p12.har.hol-type-correctness}{thm:p12-har-hol-type-correctness}{Type correctness without normalization}
\begin{proof}
If \(\Gamma\vdash M:A\), then either \(A\) is a sort or
\(\Gamma\vdash A:s\) for a sort \(s\).
Induct on the derivation. Sort and product conclusions have sort types.
Variable supplies the annotation's sort judgment, carried to its
extended context by weakening. Weakening carries the type-correctness
judgment from its subject premise; abstraction supplies the product's
sort judgment; conversion supplies the new type's sort judgment.
For application, apply the induction hypothesis to its function premise.
Its literal product type is not a sort, so it is well typed at some sort.
\cref{thm:p12-har-hol-generation} applied to that product yields its domain and codomain sort premises;
the codomain sort is the product's sort, since sorts do not beta-convert.
\Cref{thm:p12-har-pts-substitution}, with the actual argument
premise, yields \(\Gamma\vdash C[Q/x]:s_2\). This proves the remaining
case. No subject-reduction hypothesis entered \cref{thm:p12-har-hol-type-correctness}.
\end{proof}
\end{annexitem}

\begin{annexitem}{annex:p12.har.hol-context-conversion}{thm:p12-har-hol-context-conversion}{Typed context conversion}
\begin{proof}
Suppose \(\Gamma,x:B,\Theta\) is valid,
\(\Gamma\vdash B':s'\), and \(B=_\beta B'\).
Every judgment in the old context is derivable in
\(\Gamma,x:B',\Theta\), with its subject, type, and later raw
annotations unchanged. The new context is valid.
Use induction on the source derivation, split at the selected declaration.
Premises before \(x\) are unchanged, by the free-variable discipline
already established in the supplied substitution proof. At the selected
variable case derive \(x:B'\), then convert to \(x:B\): the original
validity certificate \(\Gamma\vdash B:s\) is available in the new
extension by weakening. At a later variable or weakening step transform
its annotation judgment first, making the new prefix valid, then apply
the same variable or weakening rule. Product and abstraction transform
the base and fresh extended premises and keep the same sort pair;
application transforms both premises; conversion transforms its subject
and result-type premises and keeps its raw equality. A sort case before
the selected declaration is unchanged. These cases also construct all
later annotation certificates, proving new-context validity. Reversing
the roles of \(B,B'\) proves the reverse transport when both are typed
in the prefix. This does not require subject reduction.
\end{proof}
\end{annexitem}

\begin{annexitem}{annex:p12.har.disambiguation}{thm:p12-har-disambiguation}{The complete three-stratum disambiguation construction}
\begin{proof}
Choose disjoint variable classes \(V_D,V_L,V_P\). Domain expressions and
object expressions are the following auxiliary syntactic classes:
\[
 D ::= \mathrm{Prop}\mid a\ (a\in V_D)\mid D\to D,
\]
\[
 O ::= x\ (x\in V_L)\mid\lambda x:D.O\mid OO
       \mid\Pi x:D.O\mid\Pi z:O.O,
\]
where object binders belong to \(V_L\), proof binders to \(V_P\), and
the last object's codomain contains no occurrence of its proof binder.
The notation \(D\to E\) is its actual source product with a fresh,
unused binder; it is not an additional constructor. This broad object
grammar is not a typing claim: only source-typed members will be used.
In particular its product clauses need not be well typed without the
source formation premises.

Every \(D\) is beta normal. Every beta reduct of an \(O\) is again an
\(O\), including reduction in annotations. To prove the latter, first
prove structurally that substituting an object for an object variable
preserves \(O\): domain annotations contain no object variables;
fresh binders permit induction on bodies; a proof binder has no occurrence
in an object body. A root object beta redex thus has an object contractum,
and compatible subterm contraction preserves the grammar. The classes
\(D,O,\{\mathrm{Type}\},\{\mathrm{Type}'\}\) are pairwise disjoint:
the variable classes differ; products in \(D\) end recursively in a
domain atom, whereas products in \(O\) have an object body; abstractions
and applications occur only in \(O\); the remaining sorts are neither.
Consequently no members of different such classes are beta-equal, by
their reduction closure and raw confluence. Beta-equal domain expressions
are literally equal up to their irrelevant binder names.

For each valid source context, assign its variables incrementally to the
three classes by its annotation certificates. Every source derivation
then has exactly one of these four forms:
\begin{enumerate}
\item \(\Gamma\vdash\mathrm{Type}:\mathrm{Type}'\);
\item \(\Gamma\vdash D:\mathrm{Type}\), with \(D\) in the domain grammar;
\item \(\Gamma\vdash O:D\), with \(D\) in the domain grammar and
      \(\Gamma\vdash D:\mathrm{Type}\);
\item \(\Gamma\vdash p:F\), with \(F\) in the object grammar and
      \(\Gamma\vdash F:\mathrm{Prop}\).
\end{enumerate}
Here a declaration is respectively a domain declaration \(a:\mathrm{Type}\),
an object declaration \(x:D\), or a proof declaration \(z:F\).
The third case includes formulas themselves, at the domain
\(D=\mathrm{Prop}\). The fourth case imposes no new grammar on the
proof subject; its stated formula typing is carried as part of the invariant.

Prove context classification and this four-case invariant simultaneously
by induction on annotation certificates and typing derivations. The sort
axioms give cases two and one. In a variable rule, the annotation premise
at a sort can be only case one (annotation Type), case two (annotation
domain), or case three at domain Prop (annotation formula); the new
variable consequently has its prescribed domain/object/proof class.
Weakening assigns that class and retains the invariant for its subject.
For product formation with result sort Type both domain and codomain
are domain expressions. Its binder is an object variable and so cannot
occur in its domain-grammar codomain; this is a nondependent domain arrow.
With result sort Prop, the codomain is an object expression at domain
Prop. A Type-sorted annotation makes an object binder and gives
\(\Pi x:D.F\); a Prop-sorted annotation makes a proof binder, which
cannot occur in the object codomain, and gives \(\Pi z:F.G\).
Thus both admitted cases are in the object grammar at domain Prop.

For abstraction, its well-formed product premise has one of precisely
these product forms. A domain-arrow type requires a domain annotation
and object body, so the resulting lambda belongs to the object grammar.
A propositional-product type gives case four, and that product's
Prop typing is already its premise. For application with a domain-arrow
function type, the domain and codomain are rigid domain expressions;
function and argument are objects and the application is an object at
the unchanged, nondependent codomain. With a propositional function type,
\cref{thm:p12-har-hol-generation} gives a Prop-sorted object codomain. An object argument substitutes an
object for an object variable, preserving that grammar. A proof argument
substitutes for a variable absent from the object codomain, leaving it
unchanged. The supplied substitution theorem gives the resulting formula's
Prop judgment, establishing case four. These cases exhaust applications.

For conversion, its well-typed target at a sort is either Type, a domain,
or an object formula, by the induction hypothesis. Its old type belongs
to one of the four separated type classes. Raw confluence therefore
forces the same class. A domain target and old domain are identical,
so the object invariant persists. Formula conversion retains case four
with the new target's explicit Prop judgment. Conversion to Type retains
case two; conversion from the exceptional Type' result is impossible,
since no well-typed target belongs to that class. This is the step that
rules out spurious syntax allowed by arbitrary untyped beta expansions;
neither a normalization theorem nor untyped expansion closure is used.

The induction proves more than free-variable absence: domain-expression
judgments reconstruct in the domain subcontext alone, and object judgments
reconstruct in its domain-and-object subcontext, in their original relative
order. For domains use their structural formation: a domain atom has its
axiom/projection, and a domain arrow uses Type/Type with a weakened codomain.
For objects use the same induction cases above. Object projection retains
its object declaration. Object abstraction and application retain just
their object premises. A universal-product clause retains the object
binder and its two premises. An implication-product clause first reconstructs
its two object formulas without the proof binder, then weakens the
codomain by that binder solely to apply the source Prop/Prop formation
rule. An object conversion has identical rigid domain types and can be
omitted. Every ambient proof weakening is omitted; domain/object weakening
is retained. This supplies actual smaller-context derivations rather
than appealing to a general unproved strengthening principle.

Now form \(\Gamma_D\) from all domain declarations, then
\(\Gamma_L\) from all object declarations, then \(\Gamma_P\) from all
proof declarations, keeping relative order within each. Domain declarations
have the constant annotation Type, so their reordered context is valid.
Every object annotation is a domain formed in \(\Gamma_D\); therefore
\(\Gamma_D,\Gamma_L\) is valid. Every proof annotation is a formula
reconstructed in the old domain/object subcontext; \cref{thm:p12-har-hol-contexts} carries it into
\(\Gamma_D,\Gamma_L\), so \(\Gamma_P\) can be appended. The permutation
is valid and has exactly the old declarations, hence \cref{thm:p12-har-hol-contexts} carries every
old judgment to it. All assertions of source Lemma~3.11 follow:
domains are simple arrows formed in \(\Gamma_D\); objects need no proof
declarations; a Prop/Prop product's body does not contain its proof binder.
In particular Type/Type dependence is impossible as well.
\end{proof}
\end{annexitem}

\begin{annexitem}{annex:p12.har.hol-uniqueness}{thm:p12-har-hol-uniqueness}{Uniqueness of source lambdaHOL types}
\begin{proof}
For the exact source system, \(\Gamma\vdash M:A\) and
\(\Gamma\vdash M:B\) imply \(A=_\beta B\).
Induct structurally on the common subject \(M\), using \cref{thm:p12-har-hol-generation} to remove
terminal weakening/conversion in both derivations. Variable and sort
types are fixed by their declaration or single axiom. A product's two
codomain sort premises type the same smaller body in the same literal
extended context; the induction hypothesis makes these sorts beta-equal
and hence equal. Their product result sorts are therefore equal.
For an abstraction the annotation is the same syntactic domain in both
derivations. The body induction gives equality of the two body types;
compatible products give equality of the generated abstraction types.
For an application the function induction gives equality of its two
generated product types. \cref{thm:p12-har-hol-generation}'s product injectivity gives beta equality of
their domains and bodies (with the same fresh binder). Substitution of
the common argument preserves beta equality, so the result types are
equal. Finally compose with the two generated outer type equalities.
These are every possible subject constructor. The proof needs only
generation, raw confluence, and the exact functional sort/product rules;
it assumes neither subject reduction nor checker uniqueness.
\end{proof}
\end{annexitem}

\begin{annexitem}{annex:p12.har.hol-subject-reduction}{thm:p12-har-hol-subject-reduction}{Beta subject reduction including annotations}
\begin{proof}
If \(\Gamma\vdash M:T\) and \(M\longrightarrow_\beta M'\), then
\(\Gamma\vdash M':T\). Induct structurally on \(M\), with the induction
statement quantified over every typing derivation and valid context for
that smaller subject. This avoids treating generated witnesses as if
they necessarily had smaller derivation heights. \cref{thm:p12-har-hol-generation} yields a constructor
judgment at generated type \(T_0=_\beta T\).
If the subject reduces, \cref{thm:p12-har-disambiguation} excludes its rigid domain and sort cases;
it is an object or proof, whose type \(T\) has a sort by \cref{thm:p12-har-hol-type-correctness} and \cref{thm:p12-har-disambiguation}.
Thus once the reduced constructor has type \(T_0\), conversion to
the original \(T\) is licensed. The following are the constructor cases.

For \(M=\Pi x:A.B\), a domain-annotation contraction gives
\(\Gamma\vdash A':s_1\) by the smaller-subject induction on \(A\).
\cref{thm:p12-har-hol-context-conversion} carries \(B:s_2\) to \(\Gamma,x:A'\), and the same allowed pair
rebuilds the product at \(s_2\). A body contraction uses the induction
on \(B\) under the old extension and rebuilds the same product.
For \(M=\lambda x:A.P\), \cref{thm:p12-har-hol-generation} applied also to its well-formed product
gives \(A:s_1\), \(B:s_2\). A body contraction uses the induction
on \(P:B\). An annotation contraction uses the induction on
\(A:s_1\), then \cref{thm:p12-har-hol-context-conversion} to carry both \(P:B\) and \(B:s_2\) to the
new binder context. Product formation and abstraction give the new
lambda type \(\Pi x:A'.B\), beta-equal to the old generated type.
The original product type is well formed, so conversion restores it
before any outer conversion. This explicitly preserves typed annotations.

For \(M=PQ\), a contraction in either proper child uses its
smaller-subject induction with the generated product or domain type,
then the same application rule. At a root redex
\((\lambda x:A.P)Q\), generation of the function judgment yields
\(\Gamma,x:A\vdash P:B\), a well-formed product
\(\Pi x:A.B\), and its equality to the application's product
\(\Pi x:C.E\). Product injectivity gives \(A=_\beta C\),
\(B=_\beta E\). Its formation gives \(\Gamma\vdash A:s_1\),
so convert the actual argument \(Q:C\) to \(Q:A\). Substitution
gives \(\Gamma\vdash P[Q/x]:B[Q/x]\). The generated application
type \(E[Q/x]\) has a sort by the carried typing invariants in \cref{thm:p12-har-disambiguation}:
a domain is formed at Type, and an object formula is typed at Prop.
This also covers the literal domain Prop, for which \cref{thm:p12-har-hol-type-correctness}'s weaker
``type is a sort'' alternative alone would not supply the sort judgment.
Substitution preserves \(B=_\beta E\), so convert to
\(E[Q/x]\), then to \(T\). Sorts and variables have no contraction.
These cases cover beta at every source-compatible position.
Induction on a finite path yields multi-step subject reduction.
For two paths from a term of type \(T\), raw confluence supplies a
common reduct. The endpoints retain type \(T\) by this finite-path
result, and reducing either endpoint to the common reduct preserves that
same type again. Thus the common reduct is typed, proving the stated
typed-confluence consequence as well.
\end{proof}
\end{annexitem}

\begin{annexitem}{annex:p12.har.hol-cc-transport}{thm:p12-har-hol-cc-transport}{Judgment-dependent transport to exact CC}
\begin{proof}
Let the target be precisely source CC:
sorts \(\ast,\Box\), axiom \(\ast:\Box\), and all four product
pairs \((\ast,\ast),(\ast,\Box),(\Box,\ast),(\Box,\Box)\).
Define a homomorphism \(h\) on raw pseudo terms by
\(h(\mathrm{Prop})=h(\mathrm{Type})=\ast\),
\(h(\mathrm{Type}')=\Box\), fixing variables and the product,
abstraction, and application constructors. It commutes with substitution
by structural induction and preserves each designated beta contraction,
including contractions in annotations. It is not a PTS morphism:
the source axiom Prop:Type would homomorphically become the false
target axiom \(\ast:\ast\).

For a domain expression \(D\), define \(\sigma(D)=\ast\) when its
rightmost atom is a declared domain variable and \(\sigma(D)=\Box\)
when that atom is Prop. Equivalently, \(\sigma(a)=\ast\),
\(\sigma(\mathrm{Prop})=\Box\), and
\(\sigma(D\to E)=\sigma(E)\).
Map source contexts homomorphically on annotations; in particular
\(a:\mathrm{Type}\) becomes \(a:\ast\). Map their classified
judgments as follows:
\[
\begin{array}{rcl}
 \Gamma\vdash\mathrm{Type}:\mathrm{Type}'
   &\mapsto& h\Gamma\vdash\ast:\Box,\\
 \Gamma\vdash D:\mathrm{Type}
   &\mapsto& h\Gamma\vdash hD:\sigma(D),\\
 \Gamma\vdash O:D
   &\mapsto& h\Gamma\vdash hO:hD,\\
 \Gamma\vdash p:F
   &\mapsto& h\Gamma\vdash hp:hF.
\end{array}
\]
The result-sort override is part of this typed judgment map; it is not
an assumption that the homomorphic sort image types every image term.

Prove valid target contexts and these judgment translations simultaneously
by induction on the \cref{thm:p12-har-disambiguation}-classified source derivations. Both source sort
axioms become the target axiom \(\ast:\Box\). A domain declaration
becomes a type variable licensed by that axiom. An object declaration
is licensed by \(hD:\sigma(D)\); a proof declaration is licensed
by \(hF:\ast\). Variable and weakening consequently work at the
translated sorts. A source domain product becomes a target product with
domain sort \(\sigma(D)\), codomain sort \(\sigma(E)\), and result
sort \(\sigma(E)\). All four possibilities are admitted by exact CC.
A formula product has target codomain sort \(\ast\); its domain
sort is \(\sigma(D)\) for a data annotation, or \(\ast\) for a
formula annotation. The corresponding \((\ast,\ast)\) or
\((\Box,\ast)\) pair is admitted. Abstraction uses the translated
body and exactly this translated well-formed product. Application uses
the translated product and argument; substitution commutation identifies
the target result with the mapped source result.

In a source conversion of an object/proof type, the target type has its
translated sort judgment and the raw beta equality maps by \(h\), so
target conversion is legal. A source conversion in a domain judgment
retains the literal result Type by \cref{thm:p12-har-disambiguation}'s separation, and the source domain
is unchanged; its already translated \(hD:\sigma(D)\) is retained
without an invalid attempted conversion to \(\ast\). The exceptional
top-result judgment has no conversion case, again by \cref{thm:p12-har-disambiguation}. These are all
rules, establishing transport on every actual Figure~2 derivation.

For example, with \(a:\mathrm{Type}\), the source domain
\((a\to\mathrm{Prop})\to a\) maps to
\((a\to\ast)\to a:\ast\). Its domain has target sort \(\Box\)
and its codomain has target sort \(\ast\), using exactly the CC pair
\((\Box,\ast,\ast)\). This is the mixed case missed by the printed
PRED-omega inverse described next.

It follows that strong normalization of this exact CC, if separately
supplied for its full compatible annotated beta reduction, entails strong
normalization of the exact source lambdaHOL. Indeed a hypothetical
infinite source path starts with a source-typed term (or source-typed
type); the translated starting subject is CC typed by the construction,
and every contraction maps to one contraction at the same position.
Thus it gives an infinite target contraction path, contrary to that
theorem. A source top sort appearing only as a type maps to \(\Box\)
and is already a beta-normal atom. Identifying two sort symbols erases
no lambda/application redex. This proves the transfer, not the target
CC normalization theorem itself.
\end{proof}
\end{annexitem}

\begin{annexitem}{annex:p12.har.sum-checks}{thm:p12-har-sum-checks}{Both source injections and case computations}
\begin{proof}
For \(p:A_1\), after abstracting over
\(\alpha:\mathrm{Prop}\), \(h_1:A_1\to\alpha\),
\(h_2:A_2\to\alpha\), the body \(h_1p\) has type
\(\alpha\). Thus the first injection has the encoded sum type;
the second has that type because \(h_2p:\alpha\) for
\(p:A_2\). Abstraction over each \(p\) gives its stated
function type. For \(x:A_1+A_2\), application to
\(C:\mathrm{Prop}\) and the two branch functions yields
\(xf_1f_2\) with the explicit first type argument, namely
\(xCf_1f_2:C\); abstraction gives the case function.
Expanding \(\mathrm{case}(f_1,f_2)(\mathrm{in}_1p)\)
first contracts its outer lambda, then the sum proof's three lambdas,
leaving \(f_1p\). The same four application contractions for
\(\mathrm{in}_2\) leave \(f_2p\). Type/Prop licenses the
quantification over Prop, and Prop/Prop licenses the branch arrows.
No elimination to a Type-sorted conclusion occurred.
\end{proof}
\end{annexitem}

\begin{annexitem}{annex:p12.har.inductive-checks}{thm:p12-har-inductive-checks}{Every fixed constructor branch and the generated recursive families}
\begin{proof}
For nat data recursion, the zero contractum is its premise \(f_1:C\).
At \(St\), the recursive call has type \(C\) and
\(f_2t:C\to C\), so the successor contractum has type \(C\).
For nat induction the zero contractum has type \(P0\), and
\(f_2t:Pt\to P(St)\) applies to the recursive call of type
\(Pt\), yielding \(P(St)\).

For List data recursion, Nil gives \(f_1:C\). In the Cons case,
\(d:D,t:\mathrm{List}\), the branch after those two arguments has
type \(C\to C\), and the recursive call on \(t\) has type
\(C\). For List induction the corresponding branch type is
\(Pt\to P(\mathrm{Cons}\,dt)\), so the same actual recursive
call proves the constructor's predicate. These calculations include
the original tail argument, not just its recursive value.

For Tree at a Leaf, data \(f_1d:C\) or proof
\(f_1d:P(\mathrm{Leaf}\,d)\) supplies the declared result.
At Join, each \(tn:\mathrm{Tree}\). Data recursion therefore gives
\(\mathrm{Rec}_{\mathrm{Tree}}f_1f_2(tn):C\), and abstracting
\(n\) gives the whole family \(\mathrm{nat}\to C\).
The branch \(f_2dt\) accepts exactly that family and returns \(C\).
For induction, the recursive proof at \(tn\) has type \(P(tn)\),
and the abstraction has type \(\Pi n:\mathrm{nat}.P(tn)\).
The induction branch accepts it and returns
\(P(\mathrm{Join}\,dt)\). Thus the indexed family is both
constructed and typed, rather than assumed as an unspecified witness.

The same construction handles every argument scheme
\(A_1\to\cdots\to A_k\to\mu\) in the source general
scheme. Apply its actual argument to \(u_1,\ldots,u_k\) to
obtain a \(\mu\)-term, then apply the concrete recursor and
abstract over those \(u_j\). The result has type
\(A_1\to\cdots\to A_k\to C\) for data recursion, or
\(\Pi u_1:A_1\cdots u_k:A_k.P(a_j\vec u)\) for induction.
At \(k=0\) it is simply the recursive value or proof for that
argument. Applying \(f_i:\mathcal B_i(C)\) or
\(f_i:\mathcal B_i(P)\) to all original arguments and then
these families gives exactly \(C\) or
\(P(\mathrm{constr}_i\vec a)\). Nonrecursive arguments add
no recursive hypothesis. The strict-positive absence of \(\mu\)
from the \(A_j\) is the hypothesis allowing these fixed argument
domains; a negative-domain case is not covered by this construction.

There is also a direct recursion/induction argument on a supplied
well-founded constructor tree. At a node assume the child subtrees
already have their unique recursive assignments. Their assignments
agree wherever the subtree domains overlap by uniqueness, so their
union supplies all child values; each indexed argument obtains its
value family by evaluating that union at the corresponding child root.
Apply the node's branch to its original arguments and those families
to assign its own value. The same operation with child predicate
proofs supplies its induction proof. Well-founded induction proves
existence of this assignment on each subtree and uniqueness, since
any other assignment agrees on all children and uses the same branch
at the node. For set-sized node/index families the resulting maps are
formed from their uniquely determined values, not arbitrary choices.
This constructs the fold on the given well-founded tree. It does not
prove the existence, universe closure or global normalization of every
possible formal inductive extension from a constructor name alone.

For the worked evenness term, \(d0\) computes to zero and evenness
there to \(\top\), so \(\mathrm{true}\) has the base type
\(\mathrm{even}(d0)\). At a successor, doubling computes to
\(S(S(dx))\), and the two evenness recursor steps give
\(\neg\neg\mathrm{even}(dx)\). Given
\(h:\mathrm{even}(dx)\) and \(z:\neg\mathrm{even}(dx)\),
application \(zh:\bot\) gives the double-negation introduction.
Abstracting over \(z,h,x\) yields the required step
\(\Pi x:\mathrm{nat}.\mathrm{even}(dx)\to
\mathrm{even}(d(Sx))\). The actual nat induction rule now gives
the displayed full proof term's type.

For map, its branch has type
\(D\to\mathrm{List}\to\mathrm{List}\to\mathrm{List}\):
the final argument \(h\) is the already mapped tail and
\(\mathrm{Cons}(fa)h\) is a list. Recursion has type
List-to-List, so application to the bound \(l\) returns a list.
Nil computes to Nil. At Cons, the recursor computation gives
\(\mathrm{Cons}(fa)(\mathrm{Rec}_{\mathrm{List}}
\mathrm{Nil}\,\mathrm{branch}\,t)\), which is exactly the
abbreviation-expanded \(\mathrm{Cons}(fa)(\mathrm{map}\,ft)\).
These are the claimed computations and local typing checks, not a
global beta/iota soundness or normalization proof.
\end{proof}
\end{annexitem}

\begin{annexitem}{annex:p12.har.inductive-projections}{thm:p12-har-inductive-projections}{The dependent constructor and the two proposition-only eliminators}
\begin{proof}
The first dependent-constructor branch
\(\lambda z:B.\lambda y:Qz.z\) has type
\(\Pi z:B.Qz\to B\), using the stated added Prop/Type rule.
Data recursion gives \(\pi_1:\mu\to B\), and its constructor
step on \(\mathrm{In}\,bq\) beta/iota-reduces to \(b\).
The proposed predicate \(P(x)=Q(\pi_1x)\) is well typed into
Prop. At a constructor its required branch target is
\(Q(\pi_1(\mathrm{In}\,zy))\), which converts to \(Qz\)
by that first-projection computation. Thus the second branch
\(\lambda z:B.\lambda y:Qz.y\) has the required target,
induction gives \(\pi_2:\Pi x:\mu.Q(\pi_1x)\), and its
constructor computation returns \(q\). The second projection
has not been checked before the first computation it depends on.

For primitive conjunction, \(f:\varphi\to\psi\to P\)
applied to the actual Pair arguments \(u:\varphi,v:\psi\)
has type \(P\); this is the root contractum. Choosing the first
or second projection branch returns its corresponding argument. For
Eq, the reflexive constructor is at endpoints \(x,x\), and
the branch \(f:\Pi z:D.Pzz\) applied to \(x\) has precisely
the constructor target \(Pxx\). No predicate on the equality
proof and no data-valued target is introduced by these cases. The
same fixed-parameter branch calculation validates the source parametric
List instances when their separately proposed family is admitted; it
does not form that polymorphic family inside the unchanged base system.
\end{proof}
\end{annexitem}
\begin{annexitem}{annex:p12.har.cc-normalization}{thm:p12-har-cc-normalization}{Complete compatible annotated beta normalization for exact CC}
\begin{proof}
The following is an independently worded reconstruction of the direct
two-interpretation argument in Geuvers's December~1994 Computing Science
Report~94/50, sections~2--3, printed pp.~2--10. This is a separately
identified report manifestation, not a claim of byte identity with the
1995 LNCS publication cited by the 2001 chapter. The report's CC rules
are exactly the two-sort Figure~3 rules required by \cref{thm:p12-har-hol-cc-transport}. Here the
candidate notion is the reduction-closed neutral-expansion version,
rather than the report's key-redex saturated-set version; every closure
property actually used is proved below. The argument uses ordinary sets
of finite syntax, powersets and finite iterated function sets; it needs
no inaccessible cardinal or set-theoretic model of an impredicative
universe as its own member. No CC normalization theorem is a premise.

\paragraph{(a) CC structural prerequisites, with the untyped top case.}
For the exact CC rules, reapply \cref{thm:p12-har-hol-contexts}'s validity/projection/thinning,
\cref{thm:p12-har-hol-generation}'s generation, \cref{thm:p12-har-hol-type-correctness}'s type-correctness, \cref{thm:p12-har-hol-context-conversion}'s context conversion, and
\cref{thm:p12-har-hol-uniqueness}'s structural uniqueness inductions. Their respective rule cases use
only the declared axiom, the actual admitted product pair, substitution,
and raw product injectivity; replace the two lambdaHOL axioms by
\(\ast:\Box\) and use the four exact CC pairs. In particular the
application type-correctness case still explicitly derives its result
type at the codomain's sort by substitution. In the uniqueness product
case the result is again the codomain sort, so equality of the two body
sorts gives equality of the result sorts. Nothing in these five inductions
uses the lambdaHOL-specific domain/object grammar of \cref{thm:p12-har-disambiguation}.

For precision, CC subject reduction is also supplied here rather than
assumed from \cref{thm:p12-har-hol-subject-reduction}, whose domain argument was lambdaHOL-specific.
Use structural induction on the subject, universally over its typings.
For product and lambda annotation contractions, use the smaller
annotation judgment, transport the body judgment by context conversion,
and rebuild the same admitted product or abstraction. Body contractions
use the smaller body judgment. In an application, proper-child contractions
use the generated function/argument premises. At a root lambda redex,
generation and product injectivity convert the actual argument to the
literal lambda annotation; substitution types the contractum at the
literal body instance; the generated application's result type is
explicitly well typed by the application case of type correctness, so
conversion restores it. These are the \cref{thm:p12-har-hol-subject-reduction} construction steps, but here
the existence of that result-type sort is supplied directly by CC type
correctness, not by \cref{thm:p12-har-disambiguation}.

Outer conversion is licensed for every original type other than
\(\Box\): a nonsort has its sort by type correctness, and the only
other sort \(\ast\) has type \(\Box\) by the axiom and thinning.
For original type \(\Box\), no final conversion can have that target,
since no judgment types \(\Box\); generation proves this absence by
induction on the typing rules. Strip final weakenings. A variable of
literal type \(\Box\) would require a forbidden annotation certificate;
a lambda has a literal product type; an application of literal result
\(\Box\) contradicts its explicitly derived result-type sort judgment.
Thus the subject is \(\ast\), which does not reduce, or a product whose
formation result is literally \(\Box\). Product contractions rebuild
that same sort without an illicit conversion premise. This exhausts the
exception and establishes full compatible CC beta subject reduction.
Together with raw confluence it supplies typed common reducts; no
normalization was used. Raw confluence here is \cref{thm:p12-har-pseudo-confluence} with the sort alphabet instantiated
to \(\{\ast,\Box\}\), and substitution is \cref{thm:p12-har-pts-substitution}, not a new assumed CC normalization result.

If \(\Gamma\vdash A:s\), \(\Gamma\vdash B:t\), and \(A=_\beta B\),
with \(s,t\in\{\ast,\Box\}\), typed common reducts and uniqueness
give \(s=_\beta t\), hence \(s=t\). This is sort discrimination.
Also a typable subject \(\Gamma\vdash A:T\) cannot beta-convert to
\(\Box\): since that atom is normal, confluence would make it a
reduct, and subject reduction would type it. This excludes subjects,
not the normal atom \(\Box\) counted as well typed merely because
it occurs on the right of the axiom. These facts justify classification and
semantic conversion cases below independently of normalization.

\paragraph{(b) CC classification and kind shapes.}
Partition variable names into constructor variables \(\alpha\) and
object variables \(x\), by the unique sort of each declaration's
annotation, renaming binders consistently. Write \(K,C,O\) for these
auxiliary syntax classes:
\[
\begin{aligned}
 K&::=\ast\mid\Pi x:C.K\mid\Pi\alpha:K.K,\\
 C&::=\alpha\mid\lambda x:C.C\mid\lambda\alpha:K.C
       \mid CC\mid CO\mid\Pi x:C.C\mid\Pi\alpha:K.C,\\
 O&::=x\mid\lambda x:C.O\mid\lambda\alpha:K.O\mid OO\mid OC.
\end{aligned}
\]
These grammars classify syntax, not arbitrary untyped grammar members
as typable. They are pairwise disjoint. For products, compare the body
class recursively; for lambdas likewise compare the body; for applications
compare the left child recursively; variables and the atom \(\ast\)
separate the initial cases. Matching substitution preserves all three
classes: a constructor substituted for a constructor variable, or an
object for an object variable, preserves \(K,C,O\). Prove this by
simultaneous structural induction: the selected variable is replaced
in its own class; application combines its child classes; fresh binder
renaming permits the annotation/body inductions in every binder case;
the atom is unchanged.

Every typed CC subject belongs to exactly the appropriate one of these
classes: \(M:\Box\) gives \(M\in K\); \(M:A\) with
\(A:\Box\) gives \(M\in C\); \(M:A\) with \(A:\ast\) gives
\(M\in O\). In particular an expression of type \(\ast\) is a
constructor, and a kind is an expression of type \(\Box\).
Prove the classifications simultaneously by induction on derivations,
carrying the literal class of the subject's type and the context tags.
The axiom supplies \(\ast\in K\). An annotation at \(\Box\) is a
kind, so its variable is a constructor; an annotation at \(\ast\)
is a constructor type, so its variable is an object. Weakening retains
these facts. A product with codomain sort \(\Box\) has kind body and
is in \(K\), with annotation in \(C\) or \(K\) according to its
domain sort. A product with codomain sort \(\ast\) has constructor
body and is in \(C\), with exactly the analogous two annotation cases.
For a lambda, its product-sort premise exposes the codomain sort. A
kind codomain requires a constructor body; a type codomain requires an
object body. This gives the two lambda-body classes in the grammar.

For an application, the function induction classifies its literal
product type. A kind product has kind body and constructor function;
a type product has constructor body and object function. The argument
induction gives constructor or object according to the domain's sort.
The matching-substitution lemma makes the resulting kind/type instance
remain in its class. The explicit substitution sort judgment supplies
its typing. Conversion preserves the sort of every typed old/new type
by part~(a)'s sort discrimination. The untyped old top \(\Box\) cannot
convert to a well-typed target, again by part~(a). Thus conversion retains
the induction class. This treats all primitive rule and sort cases,
including the top case, without an appeal to the report's unproved
classification citation. An expression \(\Box\) itself occurs only
as the exceptional type, never inside a typed subject.

Define a finite kind shape by
\(\kappa(\ast)=\mathbf b\),
\(\kappa(\Pi x:A.B)=\kappa(B)\) for an object binder, and
\(\kappa(\Pi\alpha:A.B)=\kappa(A)\to\kappa(B)\) for a constructor
binder. Kind syntax has no variable, lambda or application at its root;
constructor variables occur only inside its ignored type annotations.
Matching constructor/object substitution therefore leaves its shape
unchanged, by structural induction. A typed kind contraction also leaves
its shape unchanged: it cannot be a root redex; contractions in type
annotations are ignored; contractions in kind annotation/body use their
structural induction. The domain's type/kind tag cannot change under
that contraction by subject reduction and sort discrimination.
Beta-equal typed kinds consequently have identical shapes, by their
typed common reduct. These statements require no normal form for a kind.

\paragraph{(c) Candidates and the actual beta-expansion lemma.}
Use finite untyped lambda terms with one inert constant \(d\), ordinary
compatible beta contraction, and \(\mathsf{SN}\) as the set having no
infinite such path. A candidate \(R\) is a subset of this one raw-term
set satisfying: (i) \(R\subseteq\mathsf{SN}\); (ii) it is closed
under contraction; (iii) a neutral term whose every immediate reduct is
in \(R\) belongs to \(R\). Here variables, \(d\), and applications
are neutral; only a lambda is nonneutral. A redex application is neutral
and is not exempted from (iii). Candidates exist: \(\mathsf{SN}\)
itself satisfies these clauses, since any infinite path from a neutral
term must start at an immediate reduct. Every candidate contains every
variable and \(d\), since they have no reducts.

For a strongly normalizing finite term \(P\), let \(\eta(P)\) be its
greatest reduction length. This is finite: there are finitely many redex
positions at each finite term; if arbitrarily long paths existed, at
each finite branching one could choose a child admitting arbitrarily
long continuations, constructing an infinite path. Reducts strictly
lower \(\eta\). Subterms of an SN term are SN, since any subterm path
lifts through its fixed context. An inert-headed application
\(dP_1\cdots P_n\) is SN when all \(P_i\) are SN: every step lowers
\(\sum_i\eta(P_i)\). A lambda is SN exactly when its body is SN.

The controlled expansion lemma is: if \(F,N\in\mathsf{SN}\) and
\(F[N/x]\in R\), then \((\lambda x.F)N\in R\).
Induct on \(\eta(F)+\eta(N)\). Its immediate root contraction is
the given member. A contraction in \(F\) gives \(F'\); substitution
preserves that contraction, so \(F'[N/x]\in R\) by reduction closure,
and the smaller induction gives \((\lambda x.F')N\in R\).
A contraction \(N\to N'\) gives
\(F[N/x]\to^*F[N'/x]\) by reducing its finitely many substituted
copies, possibly zero, so reduction closure and the smaller induction
give \((\lambda x.F)N'\in R\). These are all immediate reducts;
neutral expansion proves the desired membership. Substitution preservation
and the finite-copy statement follow structurally on untyped syntax,
freshening binders, by the same capture-avoiding substitution identity
as in the stated substitution proof. This lemma does not assert
arbitrary beta expansion preserves SN.

Nonempty intersections of candidates are candidates, clause by clause.
The arrow
\(R\Rightarrow S=\{F:\forall N\in R,\ FN\in S\}\)
is also a candidate. First choose a fresh variable in \(R\): if
\(F\) were not SN, its path would lift into \(Fx\in S\), a
contradiction. Reduction closure follows by lifting \(F\to F'\)
into each application and using closure of \(S\). For neutral
expansion, suppose every immediate reduct of neutral \(F\) belongs
to the arrow. For \(N\in R\), induct on \(\eta(N)\).
Every immediate reduct of \(FN\) is either \(F'N\), in \(S\)
by the assumption, or \(FN'\), in \(S\) by closure of \(R\) and
the smaller induction. There is no root contraction because \(F\)
is not a lambda. Thus \(FN\in S\) by neutral expansion. This proves
the last candidate clause and finishes arrow closure.

Arrow introduction is now justified: if \(F[N/x]\in S\) for every
\(N\in R\), then \(\lambda x.F\in R\Rightarrow S\).
Take \(N=x\), after the usual fresh choice, to get \(F\in S\)
and hence its SN premise; all other \(N\) are SN by membership in
\(R\). Controlled expansion gives the required applications.
Let \(\mathrm K=\lambda p.\lambda q.p\). Applying this introduction
twice shows \(\mathrm K\in R\Rightarrow(\mathsf{SN}\Rightarrow R)\)
for every candidate \(R\). Hence
\(F\in R,\ A\in\mathsf{SN}\) implies \(\mathrm KFA\in R\).
This is the annotation-retention lemma used below, with both premises
proved rather than hidden inside an expansion assumption.

\paragraph{(d) Finite semantic domains and canonical inhabitants.}
Let \(\mathcal C\) be the set of all candidates, a subset of the
powerset of the fixed raw-term set. Interpret a finite kind shape by
\(V(\mathbf b)=\mathcal C\) and
\(V(k\to\ell)=V(\ell)^{V(k)}\), the ordinary function set.
Write \(V(A)=V(\kappa(A))\) for a source kind. Each is a set by
finite recursion, and the collection over all finite shapes is a set
by replacement over the countable finite-shape syntax. No universe
closure axiom is used. Each is nonempty with explicit canonical element:
\(c_{\mathbf b}=\mathsf{SN}\), and
\(c_{k\to\ell}\) is the constant function of value \(c_\ell\).
Kind-shape invariance gives \(V(A[P/\alpha])=V(A[t/x])=V(A)\)
and \(V(A)=V(B)\) for beta-equal typed kinds. These equalities do not
say that all kind candidates are equal.

A semantic constructor valuation \(\xi\) assigns to every declared
\(\alpha:A\) an element of \(V(A)\). These domains do not depend
on object assignments or on \(\xi\). Define a partial structural
interpretation for constructors and kinds, and an auxiliary interpretation
of the exceptional top, by
\[
\begin{aligned}
 [\ast]_\xi&=[\Box]_\xi=\mathsf{SN},&
 [\alpha]_\xi&=\xi(\alpha),\\
 [PQ]_\xi&=[P]_\xi([Q]_\xi)&&\text{for a constructor argument }Q,\\
 [Pt]_\xi&=[P]_\xi&&\text{for an object argument }t,\\
 [\lambda\alpha:A.P]_\xi&=(a\mapsto[P]_{\xi[\alpha:=a]})
      &&(a\in V(A)),\\
 [\lambda x:A.P]_\xi&=[P]_\xi&&\text{for an object binder},\\
 [\Pi x:A.B]_\xi&=[A]_\xi\Rightarrow[B]_\xi
      &&\text{for an object binder},\\
 [\Pi\alpha:A.B]_\xi&=[A]_\xi\Rightarrow
           \bigcap_{a\in V(A)}[B]_{\xi[\alpha:=a]}
      &&\text{for a constructor binder}.
\end{aligned}
\]
The product codomain \(B\) is a type or kind, so its interpretation is
a candidate, not an arbitrary semantic function. The constructor-lambda
body \(P\) is a constructor, whose interpretation may be a function.
These distinctions are supplied by part~(b). An object does not itself have
a set interpretation; it is only ignored when occurring as an argument
of a constructor. Thus \([\ast]\) is a candidate, whereas its kind
domain \(V(\ast)\) is the set of candidates; they are not identified.

\paragraph{(e) Definedness, membership, and semantic substitution.}
For every actual typed constructor \(P:A\) with \(A:\Box\), and
every semantic valuation of its context, this interpretation is defined
and \([P]_\xi\in V(A)\). For every kind \(A:\Box\), its
interpretation is a candidate. In particular every type \(T:\ast\)
has a candidate interpretation. Prove these assertions simultaneously
by structural induction on the classified constructor/kind syntax,
quantified over its source typing derivations; part~(a) generation supplies
the typing cases. A constructor variable has the prescribed valuation
member. A kind atom gives the candidate SN. A constructor lambda over
a constructor binder gives a genuine function from \(V(A)\) to
\(V(B)\), since every extended semantic valuation is valid and the
body induction supplies its value there; these are exactly the target
kind shape's function sets. A constructor lambda over an object binder
retains the body value; its target kind domain also ignores that binder.

In a constructor application to a constructor, the function induction
gives a member of \(V(A)\to V(B)\), and the argument induction
gives its input in \(V(A)\), so the semantic application is defined
and lies in \(V(B)\). The actual result kind has this same domain by
kind-shape substitution. With an object argument, both the semantic
value and the result's kind domain are unchanged. In a product with
object binder, the two proper-subterm inductions give candidate domain
and codomain, so arrow closure gives a candidate. With constructor
binder, its proper kind-domain induction gives a candidate, and every
extended valuation gives a candidate codomain. The indexing set \(V(A)\)
is nonempty by part~(d), so intersection and arrow closure give a candidate.
For a product at sort \(\ast\), this candidate is exactly membership
in \(V(\ast)=\mathcal C\). For a product at sort \(\Box\), it
is the asserted kind-candidate property. Conversion of a constructor's
kind type retains the semantic-domain membership by kind-shape equality
from part~(b); conversion does not change the recursively evaluated subject.
A target literal \(\Box\) has no conversion case, as in part~(a).
Weakening or a changed ambient valid context changes no used assignment.
These cases exhaust the syntax and its typing rules. This argument
establishes the evaluator's definedness; it does not assume it when
applying a purported undefined function.

The exact substitution equalities are
\[
 [B[P/\alpha]]_\xi=[B]_{\xi[\alpha:=[P]_\xi]},\qquad
 [B[t/x]]_\xi=[B]_\xi,
\]
for typed matching substitutions and a constructor or kind \(B\).
Prove them by simultaneous structural induction, with fresh binders.
A constructor variable gives its selected replacement or unchanged
valuation; the atom is unchanged. Constructor application uses the two
constructor substitutions and ordinary function application. An object
argument stays an object under matching substitution and remains ignored.
At a constructor lambda, the bound variable is fresh, its value ranges
over the unchanged \(V(A)\) by part~(d), and the body induction gives the
same function pointwise. At an object lambda, both sides ignore the
binder and the body induction applies. At a product the annotation and
body inductions identify the arrow factors; for a constructor binder the
indexing domain is unchanged and the equality holds at every index, so
the intersections agree. Object substitution cannot introduce a new
constructor argument or alter a variable tag; matching-substitution
classification supplies this fact. These are all clauses.

The interpretation is invariant along every typed beta contraction.
At a constructor-binder root, its semantic function application and the
first substitution equality give the contractum. At an object-binder
root, both abstraction and argument were ignored, and the second equality
gives the contractum. A kind has no root redex. Compatible steps in
constructor children use these same identities and structural induction;
object-only children are ignored. A step in a binder annotation keeps
its type/kind tag by part~(a), keeps the indexing domain by part~(d) when it is a
kind, and keeps its candidate factor by this smaller-subterm induction.
Consequently every binder clause is preserved, including product factors
and constructor-lambda function domains. Typed common reducts from part~(a)
then show \([A]_\xi=[B]_\xi\) for beta-equal well-typed types or
kinds. This is why untyped conversion is sound for this evaluator; no
normalization or untyped expansion closure of classified syntax is used.

\paragraph{(f) The annotation-retaining untyped interpretation.}
Define \(E\) structurally on typed subjects, with \(d\) inert and
\(\mathrm K=\lambda p.\lambda q.p\):
\[
\begin{aligned}
 E(\ast)&=d,& E(v)&=v,& E(MN)&=E(M)E(N),\\
 E(\lambda v:A.M)&=\mathrm K(\lambda v.E(M))E(A),&&&\\
 E(\Pi v:A.B)&=dE(A)(\lambda v.E(B)).&&&
\end{aligned}
\]
Top \(\Box\) need not be interpreted as a subject, since part~(b) shows
it never occurs in one; for an isolated type it is already normal.
All binders are freshened for substituted images. Structural induction
proves \(E(M[N/v])=E(M)[E(N)/v]\): the variable/atom/application cases
are immediate; a binder uses the annotation/body inductions and the same
fresh name. The annotation is outside its binder in both source and
target, which preserves the exact source binding scope.

A designated root contraction is simulated by three target contractions:
\[
 \mathrm K(\lambda v.E(M))E(A)E(N)
 \longrightarrow^3 E(M)[E(N)/v]=E(M[N/v]).
\]
For a designated contraction in an annotation or body, its positive
induction path lifts through the actual corresponding target context.
For a contraction in an application child it likewise lifts through
application. The product annotation and body have both been retained
in \(dE(A)(\lambda v.E(B))\); the lambda annotation remains present
as an argument of K until its administrative contraction. Thus every
source compatible contraction, at every position, has a positive finite
target path. We do not erase an annotation first and then claim its
reductions survive. Concatenating the positive paths from an infinite
source path would give an infinite target path. Hence \(E(M)\in
\mathsf{SN}\) implies full compatible source SN, including annotations.

For a finite raw substitution \(\rho\) on the context variables, let
\(E_\rho(M)=E(M)\rho\), substituting once into the original syntax.
This definition and the preceding substitution identity give
\(E_\rho(M[N/v])=E_{\rho[v:=E_\rho(N)]}(M)\) when the bound name
is fresh for the other images. In particular the image of a binder
under \(\rho\) uses its body evaluated with its own variable left
as that variable, not recursively substituted through the inserted images.

\paragraph{(g) Fundamental lemma: every sort/product/binder case.}
Call \((\xi,\rho)\) a realized valuation of \(\Gamma\) when
\(\xi\) is its semantic constructor valuation and, for every
declaration \(v:A\), \(\rho(v)\in[A]_\xi\). Every annotation is
a type or kind, so part~(e) supplies this candidate. For every actual typing
\(\Gamma\vdash M:T\) and every realized valuation,
\[
                         E_\rho(M)\in[T]_\xi.
\]
Prove this by structural induction on the subject, universally over its
typings, using generation and semantic conversion invariance. This
induction calls only proper annotation/body/application children, not
possibly larger synthesized types. A generated outer conversion replaces
\([T_0]\) by the identical \([T]\); the top case has no conversion.
Generation removes final weakenings, or restricts the valuation to the
used prefix. The constructor cases are as follows.

For the sort \(\ast:\Box\), the image is \(d\in\mathsf{SN}
=[\Box]\). A variable's image is its realized-valuation member of
the declared type candidate; conversion invariance gives its final type.
For a product \(\Pi v:A.B:s_2\), the annotation induction yields
\(E_\rho(A)\in[s_1]=\mathsf{SN}\), for either \(s_1\).
Extend the valuation by \(\rho(v)=v\). A fresh variable belongs to
every candidate. If this is a constructor binder, also choose the explicit
canonical \(c_{\kappa(A)}\in V(A)\) for \(\xi(v)\); if it is an
object binder no semantic assignment is added. Hence the body induction
gives \(E_{\rho[v:=v]}(B)\in[s_2]=\mathsf{SN}\), for either
\(s_2\). Its lambda is SN. The inert-headed application
\(dE_\rho(A)(\lambda v.E_{\rho[v:=v]}(B))\) is therefore SN,
which is membership in the product's target sort interpretation.
This covers all four product triples, including \((\Box,\Box,\Box)\),
not just the type/object product cases printed in the report's abbreviated
soundness argument.

For an abstraction over an object \(x:A\), the generated body is
\(P:B\) and its target product interpretation is
\([A]_\xi\Rightarrow[B]_\xi\). For every \(p\in[A]_\xi\),
the extended realized valuation gives by induction
\(E_{\rho[x:=p]}(P)\in[B]_\xi\); object assignments do not change
the semantic candidate of the codomain. By the E substitution identity,
these are exactly the instances of the fresh-body term
\(F=E_{\rho[x:=x]}(P)\). Arrow introduction from part~(c) gives
\(\lambda x.F\) in the required arrow. The annotation induction
gives \(E_\rho(A)\in\mathsf{SN}\); K retention therefore puts
the whole image of the annotated source abstraction in that arrow.
This holds whether \(P\) is a constructor or an object.

For an abstraction over a constructor \(\alpha:A\), the target
candidate is \([A]_\xi\Rightarrow J\), where
\(J=\bigcap_{a\in V(A)}[B]_{\xi[\alpha:=a]}\).
For every \(p\in[A]_\xi\) and every \(a\in V(A)\), extend both
valuations by \(\rho(\alpha)=p\), \(\xi(\alpha)=a\).
The body induction gives membership in the corresponding codomain
candidate. Its raw image depends only on \(\rho\), not on \(\xi\),
so the same term belongs to the entire intersection \(J\). It follows
by arrow introduction and annotation K retention that the whole lambda
image belongs to \([A]_\xi\Rightarrow J\). This argument also
covers both object and constructor bodies and supplies the SN body
premise using the fresh-variable instance; it never assumes the source
lambda was already normalizing.

For application with object argument \(N:A\), the generated function
candidate is \([A]_\xi\Rightarrow[B]_\xi\). Function and argument
inductions therefore give \(E_\rho(M)E_\rho(N)\in[B]_\xi\).
Semantic object-substitution invariance identifies this with
\([B[N/x]]_\xi\), the actual result type candidate.
For application with constructor argument \(N:A\), the function
candidate is \([A]_\xi\Rightarrow J\) with the preceding
intersection. Application of the two induction members gives membership
in \(J\). Independently part~(e) supplies \([N]_\xi\in V(A)\), so choose
that very index. The resulting membership is in
\([B]_{\xi[\alpha:=[N]_\xi]}=[B[N/\alpha]]_\xi\) by semantic
substitution. These two application cases cover both object and
constructor functions; no semantic value is assigned to an object
argument. Products, abstractions, applications, variable, and sort are
every subject constructor, so the fundamental lemma is established.

\paragraph{(h) Open-context realization and full strong normalization.}
Every valid finite CC context has a realized valuation with raw
\(\rho(v)=v\): assign to each constructor declaration
\(\alpha:A\) the explicit canonical member
\(c_{\kappa(A)}\in V(A)\). Kind-domain shapes do not depend on any
prior semantic value, so these assignments are valid simultaneously.
part~(e) says every annotation's interpretation is a candidate; its variable
belongs to that candidate by part~(c). Thus identity \(\rho\) realizes
all object and constructor declarations, including uninhabited source
types: a raw open variable is a candidate member, not a fabricated
closed proof of that type. No closure or inhabited-context assumption
has been added.

For \(\Gamma\vdash M:T\), the fundamental lemma now gives
\(E(M)\in[T]_\xi\subseteq\mathsf{SN}\). part~(f)'s positive path
simulation implies that \(M\) has no infinite full compatible beta
path. If a well-typed pseudo term occurs only on the right of a judgment,
type correctness gives a sort judgment for it, unless it is a sort;
\(\ast\) has its axiom judgment and \(\Box\) is a normal atom.
Hence the same conclusion covers both subjects and types in the
source's definition of well typed. The proof treats arbitrary well-formed
open contexts, all four product triples, reductions beneath binders,
and contractions inside annotations/types. It assumes neither CC SN,
closed-term normalization, weak-head termination, checker completeness,
nor an impredicative set universe. The CC result has no normalization premise. Its separate application to source lambdaHOL is \cref{thm:p12-har-hol-normalization}; definitions, delta, recursors and later inductive extensions remain outside these conclusions.
\end{proof}
\end{annexitem}

\begin{annexitem}{annex:p12.har.hol-normalization}{thm:p12-har-hol-normalization}{The complete exact-system normalization transfer}
\begin{proof}
For a source subject \(\Gamma\vdash M:T\),
\cref{thm:p12-har-hol-cc-transport} constructs an exact CC typing
judgment for \(hM\), with the stated result-sort override when
\(T=\mathrm{Type}\). Hence \cref{thm:p12-har-cc-normalization}
makes \(hM\) strongly normalizing. Each designated source beta
contraction, including a contraction in an annotation or product body,
maps homomorphically to a CC beta contraction at that position.
An infinite source path would therefore give an infinite CC path,
contradicting that result. No inverse PRED-omega map is used.

For a term occurring only as a type on the right of a judgment,
type correctness gives its own sort judgment unless it is a sort
itself. Apply the preceding subject argument to the former case.
Prop and Type have their source axiom judgments, and the untyped
top sort Type' is a normal atom; consequently all the source's
well-typed pseudo terms are covered. The rule/step comparison uses
exactly the declared beta-only base; it neither adds delta/iota nor
transfers normalization to any inductive, strong-sum or warned
universe extension.
\end{proof}
\end{annexitem}

\begin{annexitem}{annex:p12.har.development-objects}{thm:p12-har-development-objects}{The source non-or object and the DN-bearing existential equivalence}
\begin{proof}
For \(p:\neg(a\lor b)\), the left injection maps
\(q:a\) to \(a\lor b\), so \(p(\mathrm{or\_inl}\,a b q)\)
has type FF. Abstracting \(q\) gives \(\neg a\).
The right injection gives \(\neg b\) in the same way.
Applying the declared \(\mathrm{and\_in}\) to these two terms
gives \(\neg a\land\neg b\), and abstraction over \(p\)
gives the source non-or object. Its logical parameters \(a,b\)
are then abstracted as in the historical displayed term.

Put \(X=\exists x:A.Px\) and
\(Y=\neg\Pi x:A.\neg Px\). Given \(H:X\) and
\(H_0:\Pi x:A.\neg Px\), existential elimination into FF uses
branch \(\lambda x:A.\lambda H_1:Px.H_0xH_1\), yielding FF.
Abstracting \(H_0\) gives \(Y\), hence a term \(X\to Y\).

For the reverse take \(H:Y\) and assume \(H_0:\neg X\).
For each \(x:A,H_1:Px\), the introduction
\(\mathrm{ex\_intro}\,A P xH_1:X\) is contradicted by \(H_0\).
Thus \(\lambda x:A.\lambda H_1:Px.
H_0(\mathrm{ex\_intro}\,A P xH_1)\) has type
\(\Pi x:A.\neg Px\). Applying \(H\) gives FF, and
abstracting \(H_0\) gives \(\neg\neg X\).
The literal source axiom \(\mathrm{DN}\,X\) now gives \(X\).
Abstracting \(H\) yields \(Y\to X\). The declared conjunction
constructor pairs these two implications into the stated equivalence.
The forward direction used no DN; the reverse used that explicit axiom
and does not become an axiom-free constructive proof by type erasure.
\end{proof}
\end{annexitem}
